% Aula 03 --- as duas formas da fórmula de k dobras (aula x ISLP).
Quem estuda validação cruzada pelas notas de aula e pelo livro encontra duas fórmulas
diferentes para a mesma quantia, e é natural desconfiar de que uma delas esteja
errada. Não está: elas só diferem na unidade que cada uma pondera igualmente.

As notas de aula definem
\[
  \riscohat_{k\text{-CV}}
   = \frac1n\sum_{j=1}^{k}\sum_{i\in L_j}\big(Y_i-g_{-j}(\X_i)\big)^2 ,
\]
e o livro define $\mathrm{CV}_{(k)}=\frac1k\sum_{j=1}^{k}\mathrm{EQM}_j$, com
$\mathrm{EQM}_j=\frac{1}{\abs{L_j}}\sum_{i\in L_j}(Y_i-g_{-j}(\X_i))^2$.
\begin{itens}
  \item \pts{1,2} Reescreva a fórmula das notas de modo a fazer os
        $\mathrm{EQM}_j$ aparecerem, mostrando que
        \[
          \riscohat_{k\text{-CV}}=\sum_{j=1}^{k}\frac{\abs{L_j}}{n}\,\mathrm{EQM}_j .
        \]
  \item \pts{0,8} Olhando para os pesos que você obteve, diga em que condição as
        duas definições coincidem --- e se coincidem exatamente ou apenas
        aproximadamente. Em seguida, nomeie a unidade que cada fórmula pondera
        igualmente.
  \item \pts{0,7} Com $n=50$ e $k=7$, as dobras não podem ter todas o mesmo
        tamanho; suponha que tenham tamanhos $8,8,7,7,7,7,6$. Confirme que somam $50$
        e explique por que, neste caso, os dois números diferem.
  \item \pts{1,3} Suponha que a dobra de $6$ observações tenha $\mathrm{EQM}=10$ e
        todas as outras tenham $\mathrm{EQM}=1$. Calcule os dois estimadores, diga
        qual deles dá mais peso à dobra atípica, e qual você reportaria. Justifique
        a escolha.
\end{itens}

\begin{solucao}
\textbf{(a)} A chave é reconhecer a soma interna. Pela definição de $\mathrm{EQM}_j$,
basta multiplicar dos dois lados por $\abs{L_j}$:
\[
  \sum_{i\in L_j}\big(Y_i-g_{-j}(\X_i)\big)^2=\abs{L_j}\,\mathrm{EQM}_j .
\]
Substituindo na fórmula das notas,
\[
  \riscohat_{k\text{-CV}}
   =\frac1n\sum_{j=1}^{k}\abs{L_j}\,\mathrm{EQM}_j
   =\sum_{j=1}^{k}\frac{\abs{L_j}}{n}\,\mathrm{EQM}_j .
\]
São os \emph{mesmos} $\mathrm{EQM}_j$ do livro, combinados com pesos
$w_j=\abs{L_j}/n$. E esses pesos somam $1$, porque as dobras particionam as $n$
observações: $\sum_j\abs{L_j}=n$. As duas fórmulas são, portanto, médias ponderadas dos
mesmos $k$ números, e só diferem nos pesos: $\abs{L_j}/n$ numa, $1/k$ na outra.

\smallskip
\textbf{(b)} Os pesos coincidem em todas as dobras se e só se $\abs{L_j}/n=1/k$ para todo
$j$, isto é, se todas as dobras têm tamanho $n/k$ --- o que só é possível quando $k$
divide $n$. Nesse caso as duas fórmulas são a mesma soma dos mesmos números com os
mesmos pesos: coincidem \textbf{exatamente}, e não aproximadamente.

A unidade que cada uma pondera igualmente: as notas dão o mesmo peso a cada
\textbf{observação} --- todo resíduo de validação entra com peso $1/n$ ---; o livro dá o
mesmo peso a cada \textbf{dobra}. Visto do lado das observações, o livro dá a cada
resíduo da dobra $j$ o peso $1/(k\abs{L_j})$.

\smallskip
\textbf{(c)} $8+8+7+7+7+7+6=50$.

Os pesos das notas são $8/50$ (duas vezes), $7/50$ (quatro vezes) e $6/50$; os do
livro, $1/7$ para cada dobra. Como os tamanhos diferem, os pesos diferem, e as duas
médias só coincidem por acaso --- quando os $\mathrm{EQM}_j$ são todos iguais, por
exemplo.

Olhando por observação fica mais claro. No livro, cada resíduo recebe
\[
  \frac{1}{7\cdot8}=\frac1{56},\qquad
  \frac{1}{7\cdot7}=\frac1{49}\qquad\text{ou}\qquad
  \frac{1}{7\cdot6}=\frac1{42},
\]
conforme a dobra em que caiu, contra $1/50$ para todos nas notas. Uma observação da
dobra de $6$ pesa $4/3$ de uma da dobra de $8$, só porque divide a sua dobra com menos
gente.

\smallskip
\textbf{(d)} Pelas notas, a dobra atípica entra com peso $6/50$ e as outras $44$
observações com $44/50$:
\[
  \riscohat_{k\text{-CV}}=\frac{6}{50}\times10+\frac{44}{50}\times1=\frac{104}{50}=\mathbf{2{,}08}.
\]
Pelo livro, a dobra atípica é uma das sete, e só:
\[
  \mathrm{CV}_{(7)}=\frac17\big(10+6\times1\big)=\frac{16}{7}\approx\mathbf{2{,}286}.
\]
O livro dá mais peso à dobra atípica: $1/7\approx0{,}143$ contra $6/50=0{,}12$. Ela é a
menor dobra e, ainda assim, vale um sétimo do total.

\emph{Qual reportar?} Para \emph{estimar o risco}, a das notas tem a seu favor um
argumento limpo. O risco é a esperança da perda numa observação nova, e cada resíduo de
validação $\ell_i$ é uma medida dessa perda. Se essas medidas tivessem a mesma média e a
mesma variância $s^2$, e fossem não correlacionadas --- uma idealização, mas é ela que
orienta ---, toda combinação $\sum_iw_i\ell_i$ com pesos somando $1$ teria a mesma
média, e variância $s^2\sum_iw_i^2$. Pela desigualdade de Cauchy--Schwarz,
\[
  1=\Big(\sum_{i=1}^n w_i\Big)^2\le n\sum_{i=1}^n w_i^2 ,
\]
com igualdade só nos pesos iguais, $w_i=1/n$. Pesar as observações igualmente é o que
menos desperdiça informação; dar mais peso a quem caiu numa dobra pequena é deixar o
acaso do corte decidir.

A do livro também se defende. Ela trata cada dobra como uma repetição do experimento
"treinar e validar", e é a escala natural quando o que se quer é comparar as dobras
entre si: é dos $k$ números $\mathrm{EQM}_j$ que sai o erro-padrão da regra de um
erro-padrão. E é a que o \texttt{scikit-learn} calcula ---
\texttt{cross\_val\_score(\dots).mean()} é média não ponderada ---, então o essencial é
\emph{saber qual das duas o seu código calcula} antes de comparar o seu número com o de
outra pessoa.

E a recomendação das notas encerra a discussão pela raiz: com $n$ pequeno, escolha um
$k$ que divida $n$. Aí as duas coincidem, e a pergunta não se coloca.
\end{solucao}
